\subsection{Verstärkungseinstellung mit schaltbaren Widerständen}
% Alt: Aufgabe 13
\Aufgabe{
    \label{Aufg13}
    Um die Verstärkung eines OPVs einzustellen, können mithilfe von Schaltern verschiedene
    Widerstände am OPV zu- und abgeschaltet werden. Berechnen Sie die untenstehenden OP-Schaltungen
    in Abhängigkeit der Schalterstellungen S!
    \begin{itemize}
        \item[ \bf a)] Für die folgende Schaltung mit vier möglichen Schalterstellungen
    \end{itemize}

    \begin{figure}[H]
        \centering
        \begin{tikzpicture}
    \ctikzset{tripoles/en amp/input height=-0.45}

    \draw (0,0) 
        node[en amp, anchor=center] (opamp) {};
    \draw (opamp.out) 
        to[short, -*] (2,0)
        to[short, -o] (3,0)node[above]{$U_\mathrm{a}$};
    \draw (2,0)
        to[R,l=$R_\mathrm{f}$, -*] (2,-2)
        to[R,l=$R_\mathrm{n1}$](2,-4)
        to[R,l=$R_\mathrm{n2}$](2,-6)node[ground]{};
    \draw (opamp.+)
        to[short, -o] (-3,0.5) node[above] {$U_\mathrm{e}$};
    \draw (opamp.-)  -- (opamp.- -| -2, 0)
        to[short, -*] (-2,-2) -- (-2,-4)
        to[normal open switch,l=$S2$, -*] (2,-4);
    \draw (-2,-2)
        to[normal open switch, l=$S1$] (2,-2);
\end{tikzpicture}\alttext{Das Bild zeigt eine Schaltung mit einem Operationsverstärker. Der Eingang \(U_\mathrm{E}\) ist am positiven Eingang des Operationsverstärkers angeschlossen. Vom Ausgang \(U_\mathrm{A}\) des Operationsverstärkers führt ein Widerstand \(R_\mathrm{f}\) nach unten, an dessen Verbindung zwei Schalter \(S1\) und \(S2\) parallel geschaltet sind. Hinter dem Schalter \(S1\) befindet sich der Widerstand \(R_{\mathrm{n1}}\), hinter dem Schalter \(S2\) der Widerstand \(R_{\mathrm{n2}}\), der mit Masse verbunden ist. Die Verbindung zwischen den Schaltern und den Widerständen führt zurück zum negativen Eingang des Operationsverstärkers.}
        \label{fig:FigOPVschaltbar1}
    \end{figure}

    \begin{itemize}
        \item[ \bf b)] Für die folgende Schaltung mit zwei möglichen Schalterstellungen
    \end{itemize}

    \begin{figure}[H]
        \centering
        \begin{tikzpicture}
\draw
(0,0) node[en amp,  anchor=center] (opamp) {};

\draw (opamp.-) 
    to[short, -*] (-1.2,1);
\draw (-4,1)node[above] {$U_\mathrm{e}$}  to[ R, l=$R_\mathrm{n}$, o-]
    (-1.2,1);
\draw (-1.2,1) -- (-1.2,2.5)
    to[R, l=$R_\mathrm{f1}$, -*] (1,2.5)
    to[R, l=$R_\mathrm{f2}$, -*] (4,2.5)
    to[short, -*] (4,0);
\draw (opamp.out) to[short, -o] (5,0) node[above]{$U_\mathrm{a}$};
\draw (1,2.5) -- (1,4) 
    to[normal open switch, l=$S1$] (4,4)
    -- (4,2.5);
\draw (opamp.+) -- (-1.2,-2)node[ground]{};

\end{tikzpicture}
\alttext{Das Bild zeigt ein elektrisches Schaltbild mit einem Operationsverstärker. Links ist der Eingang \(U_\mathrm{E}\), der über einen Widerstand \(R_\mathrm{n}\) mit dem invertierenden Eingang des Operationsverstärkers verbunden ist. Der nicht-invertierende Eingang des Operationsverstärkers ist mit Masse verbunden. Vom Ausgang des Operationsverstärkers führt eine Verbindung zu zwei parallel geschalteten Widerständen \(R_{\mathrm{f1}}\) und \(R_{\mathrm{f2}}\), die durch einen Schalter \(S1\) verbunden sind. Der Ausgang wird als \(U_\mathrm{A}\) bezeichnet.}
        \label{fig:FigOPVschaltbar2}
    \end{figure}

}

\Loesung{
    \begin{itemize}
        \item[a)] Nicht-invertierender Verstärker

              \textbf{1. Fall (S1 geschlossen, S2 offen):}
              In diesem Zustand bilden die Widerstände $R_\mathrm{n1}$ und $R_\mathrm{n2}$ eine Serienschaltung. Die Gesamtverstärkung ergibt sich somit zu:
              \[
                  V = 1 + \frac{R_\mathrm{f}}{R_\mathrm{n1} + R_\mathrm{n2}}
              \]

              \textbf{2. Fall (S1 offen, S2 geschlossen):}
              In diesem Fall sind der Widerstand $R_\mathrm{f}$ und $R_\mathrm{n1}$ seriell geschaltet. Diese Serienschaltung bildet zusammen mit dem Widerstand $R_\mathrm{n2}$ die Rückkopplungsschleife. Die Verstärkung ist daher:
              \[
                  V = 1 + \frac{R_\mathrm{f} + R_\mathrm{n1}}{R_\mathrm{n2}}
              \]

              \textbf{3. Fall (S1 und S2 geschlossen):}
              Mit beiden Schaltern geschlossen wird der Widerstand $R_\mathrm{n1}$ überbrückt und somit nicht berücksichtigt. Die Verstärkung ergibt sich durch:
              \[
                  V = 1 + \frac{R_\mathrm{f}}{R_\mathrm{n2}}
              \]

              \textbf{4. Fall (S1 und S2 offen):}
              Wenn beide Schalter offen sind, fehlt der Rückkopplungspfad vollständig, womit die Verstärkung theoretisch unendlich groß ist:
              \[
                  V = \infty
              \]

        \item[b)] Invertierender Verstärker \\
              \textbf{1. Zustand (S1 offen):}
              Sind $R_\mathrm{f1}$ und $R_\mathrm{f2}$ seriell geschaltet, ergibt sich eine größere Rückkopplung und somit eine höhere Verstärkung. Die Verstärkung errechnet sich wie folgt:
              \[
                  V = -\frac{R_\mathrm{f1} + R_\mathrm{f2}}{R_\mathrm{n}}
              \]

              \textbf{2. Zustand (S1 geschlossen, S2 offen):}
              In diesem Zustand wird der Widerstand $R_\mathrm{f2}$ durch den geschlossenen Schalter kurzgeschlossen und nicht berücksichtigt. Die Verstärkung reduziert sich damit auf:
              \[
                  V = -\frac{R_\mathrm{f1}}{R_\mathrm{n}}
              \]
    \end{itemize}

    Siehe: Abschnitt \ref{fig: Verschaltung Verstaerker} und Tabelle \ref{tab:Grundschaltungen}
}